\documentclass[10pt,a4paper]{amsart}
\usepackage[margin=19mm]{geometry}
\usepackage{amsmath,amssymb,mathtools}
\usepackage[scr=rsfs,cal=euler]{mathalfa}
\usepackage{xfrac}
\usepackage[unicode,hidelinks,bookmarks=false]{hyperref}
\newcommand{\ie}{i.e.\ }
\newcommand{\cf}{cf.\ }
\newcommand{\resp}{respectively\ }
\newcommand{\Subsection}{\subsection}
\providecommand{\nobreakdash}{\nobreak\hbox{-}}
\setlength{\emergencystretch}{2em}
\multlinegap0pt
\usepackage{mathtools}
\usepackage{amssymb}
\newtheorem{lemma}{Lemma}
\newcommand{\R}{\mathbb R}
\newcommand{\cM}{\mathcal M}
\newcommand{\cZ}{\mathcal Z}
\newcommand{\dd}{\mathrm d}
\newcommand{\dist}{\operatorname{dist}}
\title{\LARGE A Metric with Positive Sectional Curvature on $S^2\times S^3$}
\author{Shengtao Guo, Ethan X. Fang, Junwei Lu}
\date{Source: arXiv:2608.22133v2 (CC BY 4.0)}
\begin{document}
\maketitle

\noindent\emph{Source and licence.} \LARGE A Metric with Positive Sectional Curvature on $S^2\times S^3$. Version arXiv:2608.22133v2 (CC BY 4.0), distributed by arXiv under the Creative Commons Attribution 4.0 International licence, http://creativecommons.org/licenses/by/4.0/, as recorded in the arXiv licence metadata for this exact version and shown on the abstract page https://arxiv.org/abs/2608.22133v2. Full licence notice: \texttt{licenses/arxiv-2608.22133-CC-BY-4.0.txt}.\par\smallskip

\noindent\textit{Editorial scope.} Counted unit: the article's lemma lem:first-order-cancellation (source0.tex lines 613--622). The article's complete original proof of it is delivered with it (source0.tex lines 624--709, one \texttt{proof} environment of 86 lines). No mathematics is written, repaired or re-derived: the statement and the proof are the article's own lines, in the article's own notation, and no retained line is substituted or re-flowed.  No further source-local context is needed: every cross-reference of every retained line resolves inside the two delivered spans.  Nothing was omitted from any retained span, so the package carries no omission record.  No retained line contains a citation, so the wrapper carries no bibliography and no bibliography entry.  No retained line contains a figure environment, an image-inclusion command or drawing code, so no image is delivered and the package declares no asset dependency.  Editorial formatting only, in addition: an AMS \texttt{amsart} preamble that restates the article's own declarations and macros used by the retained lines, namely the environments \texttt{lemma} on separate counters, together with the standard packages those lines need.  The retained environments print in this document as Lemma 1 in order of appearance, whereas the article prints the same items with the numbering of its own full document.  The complete native source of the article as distributed in the arXiv e-print for this exact version is bundled unchanged as \texttt{sources/208306/source0.tex}.
\begin{lemma}                                                                       \label{lem:first-order-cancellation}
On $\mathscr F$, one has $\mathfrak a_0=0$, and
$\left.\partial_t\mathfrak a_t\right|_{t=0}=0$ on
$\rho^{-1}(\cM)$.  Moreover, $\mathfrak a_t=\mathfrak b_t=0$ on
$\rho^{-1}(\cZ)$.  Consequently,
\begin{equation}                                                                    \label{eq:a-b-distance-estimates}
 |\mathfrak a_t|\leq C\bigl(t\dist(\Pi,\cM)+t^2\dist(\Pi,\cZ)\bigr)
 \;\text{and}\; |\mathfrak b_t|\leq C\dist(\Pi,\cZ).
\end{equation}
\end{lemma}

\begin{proof}
Let $\bar g$ be the unit product metric, so that
$g_0=\frac12\bar g$.  Set $J_i x=p_i\times x$,
$J=J_1\oplus J_2$, and $\beta(x_1,x_2)=J_1x_1+J_2x_2$.
Throughout this first-order calculation, $\nabla=\nabla^{g_0}$, and an
overdot denotes $\left.\partial_t\right|_{t=0}$.
Then $\Omega_0(X,Y)=g_0(JX,Y)$.  Differentiating
$g_t(X,Y)=g_0(C_tX,Y)$ and the definition of $\Omega_t$ at $t=0$ gives
$\dot g(X,Y)=-\tfrac12\langle\beta(X),\beta(Y)\rangle$ and
$\dot\Omega=\tfrac12\dd\alpha$.
Terms arising from the variation of the frame do not contribute because they
are contracted with $\nabla^{g_0}\Omega_0=0$.

Choose extensions that are parallel at the point with respect to the product
connection.  Then
$(\nabla_Z\beta)(X)=z_1\times x_1+z_2\times x_2$, and hence
\begin{equation}\label{eq:metric-variation-derivative}
 (\nabla_Z\dot g)(X,Y)=-\frac12\{\langle(\nabla_Z\beta)(X),\beta(Y)\rangle+\langle\beta(X),(\nabla_Z\beta)(Y)\rangle\}.
\end{equation}
For $E=(e,0)$ and $F=(0,f)$, direct differentiation of $\alpha$ gives
$\dd\alpha(E,F)
=\langle e\times p_2,f\rangle+\langle p_1\times f,e\rangle$.
Writing $D$ for the flat connection on $\R^3$, use
$D_e e=-|e|^2p_1$ and $D_f f=-|f|^2p_2$ to obtain
\begin{equation}                                                                    \label{eq:omega-variations}
 (\nabla_E\dot\Omega)(E,F)=-\frac12|e|^2\langle p_1\times p_2,f\rangle
 \;\text{and}\; (\nabla_F\dot\Omega)(F,E)=\frac12|f|^2\langle p_1\times p_2,e\rangle.
\end{equation}
Let $\dot{\nabla}=\left.\partial_t\nabla^{g_t}\right|_{t=0}$.  The
Levi-Civita variation
formula is
\begin{equation}                                                                    \label{eq:connection-variation}
 2g_0(\dot{\nabla}_X Y,Z)
 =(\nabla_X\dot g)(Y,Z)+(\nabla_Y\dot g)(X,Z)
 -(\nabla_Z\dot g)(X,Y).
\end{equation}
Substitution of \eqref{eq:metric-variation-derivative} into
\eqref{eq:connection-variation}, first with $Z=JF$ and then with
$Z=JE$, yields
\begin{equation}                                                                    \label{eq:connection-variations}
 \begin{aligned}
 \Omega_0(\dot{\nabla}_E E,F)&=0
 &\text{and } \Omega_0(E,\dot{\nabla}_E F)&=-\frac12|e|^2\langle p_1\times p_2,f\rangle,\\
 \Omega_0(\dot{\nabla}_F F,E)&=0
 &\text{and } \Omega_0(F,\dot{\nabla}_F E)&=\frac12|f|^2\langle p_1\times p_2,e\rangle.
 \end{aligned}
\end{equation}
For instance,
$2g_0(\dot{\nabla}_E F,J E)=-|e|^2\langle p_1,J_2f\rangle
=-|e|^2\langle p_1\times p_2,f\rangle$, which is the second identity in
\eqref{eq:connection-variations}.

Let $\mathcal V_\Omega$ denote the first variation of
$\nabla^{g_t}\Omega_t$ at $t=0$; explicitly,
\[
 \mathcal V_\Omega(Z;X,Y)=(\nabla_Z\dot\Omega)(X,Y)
 -\Omega_0(\dot{\nabla}_Z X,Y)-\Omega_0(X,\dot{\nabla}_Z Y).
\]
Equations \eqref{eq:omega-variations} and \eqref{eq:connection-variations} give
$\mathcal V_\Omega(E;E,F)=\mathcal V_\Omega(F;F,E)=0$.
The tensor $\mathcal V_\Omega$ is linear in its first slot and alternating in the last two.
If $X=aE+bF$ and $Y=-bE+aF$ form an oriented orthonormal rotation
of the frame $E,F$, then $X\wedge Y=E\wedge F$, and therefore
$\mathcal V_\Omega(X;X,Y)=\mathcal V_\Omega(aE+bF;E,F)
=a\mathcal V_\Omega(E;E,F)-b\mathcal V_\Omega(F;F,E)=0$.
The calculation also applies when $p_1=\pm p_2$.  This proves the first two
assertions of the lemma.

On $T_b$, both $\omega_i$ vanish because each factor contributes only
one tangent direction.  Moreover, $p_1\times p_2$ is parallel to $b$,
while every tangent vector to $T_b$ is perpendicular to $b$.  Hence
$\alpha|_{T_b}=0$ and $\Omega_t|_{T_b}=0$.  Since $T_b$ is totally geodesic,
this implies $\mathfrak a_t=\mathfrak b_t=0$ on $\rho^{-1}(\cZ)$.

Finally, write
$\mathfrak a_t=t\mathfrak a^{(1)}+t^2\mathfrak a^{(2)}_t$.
The vanishing of $\mathfrak a^{(1)}$ on $\rho^{-1}(\cM)$ and smoothness in a
tubular neighborhood give
$|\mathfrak a^{(1)}|\leq C\dist(\rho(\,\cdot\,),\cM)$.  Both
$\mathfrak a_t$ and $\mathfrak a^{(1)}$ vanish on
$\rho^{-1}(\cZ)$, so
$|\mathfrak a^{(2)}_t|\leq C\dist(\rho(\,\cdot\,),\cZ)$.  This proves the
first estimate in \eqref{eq:a-b-distance-estimates}; the second follows from
the vanishing of $\mathfrak b_t$ on the smooth compact submanifold
$\rho^{-1}(\cZ)$.
\end{proof}

\end{document}
